% ============================================================
% Preamble — Bocconi MSc thesis format (Guide to the University §10.3)
%   A4, left/right margins 2.5 cm, 26–30 lines per page, 12 pt body,
%   numbered pages, no title page/abstract in body, no personal info.
% ============================================================

% ---------- encoding & fonts ----------
\usepackage[T1]{fontenc}
\usepackage[utf8]{inputenc}
\usepackage{lmodern}
\usepackage{microtype}

% ---------- geometry: 2.5 cm side margins, ~29 text lines/page ----------
\usepackage[a4paper,hmargin=2.5cm,textheight=630pt,vcentering]{geometry}
\usepackage{setspace}
\setstretch{1.5}   % 12pt x 1.5 -> ~29 lines per page (within the 26-30 rule)

% ---------- mathematics ----------
\usepackage{amsmath,amssymb,amsthm,mathtools,bm}
\allowdisplaybreaks[2]

% ---------- graphics, tables, captions ----------
\usepackage{graphicx}
% Figures are shared with the other three documents rather than duplicated, so
% \includegraphics takes a bare basename and this is the one place the location
% is written down. The chapters used to carry a literal `figures/' prefix each.
\graphicspath{{../shared/figures/}}
\usepackage{booktabs}
\usepackage[font=small,labelfont=bf,margin=1em]{caption}
\usepackage{subcaption}
\usepackage{capt-of}          % \captionof, used by the shared generated tables
\usepackage{xcolor}

% ---------- code listings ----------
\usepackage{listings}
\definecolor{codebg}{RGB}{247,247,245}
\definecolor{codekw}{RGB}{26,58,92}
\definecolor{codecm}{RGB}{110,110,110}
\definecolor{codestr}{RGB}{140,60,20}
\lstset{
  language=Python,
  basicstyle=\ttfamily\footnotesize\linespread{1.05}\selectfont,
  keywordstyle=\color{codekw}\bfseries,
  commentstyle=\color{codecm}\itshape,
  stringstyle=\color{codestr},
  backgroundcolor=\color{codebg},
  frame=single, framerule=0.3pt, rulecolor=\color{black!25},
  numbers=left, numberstyle=\tiny\color{black!45}, numbersep=7pt,
  breaklines=true, showstringspaces=false,
  aboveskip=1em, belowskip=1em, columns=fullflexible, keepspaces=true
}

% ---------- toolbox environment (step-by-step derivations) ----------
\usepackage{etoolbox}
\usepackage[most]{tcolorbox}
\newcounter{toolbox}[chapter]
\renewcommand{\thetoolbox}{\thechapter.\arabic{toolbox}}
% optional first argument: a label (placed via `phantom' so it survives
% breakable boxes); mandatory second argument: the title.
\newtcolorbox{toolbox}[2][]{%
  breakable, enhanced,
  colback=black!3, colframe=black!55,
  boxrule=0.4pt, arc=1.2mm, left=2mm, right=2mm, top=1.5mm, bottom=1.5mm,
  phantom={\refstepcounter{toolbox}\ifblank{#1}{}{\label{#1}}},
  title={\small\bfseries Toolbox~\thetoolbox\ --- #2},
  fonttitle=\normalfont, coltitle=white, colbacktitle=black!55,
  before skip=1.2em, after skip=1.2em
}

% ---------- theorem environments ----------
\theoremstyle{plain}
\newtheorem{theorem}{Theorem}[chapter]
\newtheorem{proposition}[theorem]{Proposition}
\newtheorem{lemma}[theorem]{Lemma}
\newtheorem{corollary}[theorem]{Corollary}
\theoremstyle{definition}
\newtheorem{definition}[theorem]{Definition}
\newtheorem{example}[theorem]{Example}
\theoremstyle{remark}
\newtheorem{remark}[theorem]{Remark}

% ---------- citations: author-date (Bocconi writing guide, Attachment 2) ----------
\usepackage[round,authoryear]{natbib}
\setlength{\bibsep}{0.4em}

% ---------- headers / page numbers ----------
\usepackage{fancyhdr}
\pagestyle{fancy}
\fancyhf{}
\fancyhead[L]{\small\itshape\nouppercase{\leftmark}}
\fancyhead[R]{\small\thepage}
\renewcommand{\headrulewidth}{0.3pt}
\fancypagestyle{plain}{\fancyhf{}\fancyfoot[C]{\small\thepage}%
  \renewcommand{\headrulewidth}{0pt}}

% ---------- hyperlinks ----------
\usepackage[hidelinks]{hyperref}
\urlstyle{same}

% ---------- notation macros (consistent with the research notes) ----------
% Same rendering as shared/notation.tex, which the thesis does not \input
% (it carries its own notation block). Kept in step with it deliberately: a
% chapter that cites a file by name should look the same in every document.
\providecommand{\codefile}[1]{\texttt{\small #1}}
% The thesis writes the training-set size N where the paper writes M. Shared
% generated sections use \nseq, so alias it here rather than forking the
% generator: the numbers must be identical in both documents, and a second copy
% of the file is exactly how they would stop being.
\providecommand{\nseq}{N}
% Same reason, for the autoregressive coefficient: shared/notation.tex sets
% \corr to \rho, and the shared non-Markov table is written through it.
\providecommand{\corr}{\alpha}
% The innovation variance. Chapters 5 and Appendix B already wrote it
% \sigma_\eta^2; Chapters 6, 8 and 9 wrote it q, which is also the geometric
% decay rate q(\alpha,t) of Chapter 5. One macro, so the two cannot drift apart
% again and so the collision with q stays fixed.
\newcommand{\ivar}{\sigma_\eta^2}
% Where THIS document keeps its provenance discussion. Shared generated
% sections cite it through this macro instead of a literal label, because the
% paper and the thesis file that material under different appendix names and a
% hard \ref resolves in one and prints "??" in the other.
\providecommand{\provenanceappendix}{Appendix~\ref{app:repro}}
% The companion compendium holds the background derivations that the thesis
% cites but does not carry, and the development record. Written once here so
% that the phrase is identical everywhere it appears.
\newcommand{\compendium}{the companion compendium}
\newcommand{\R}{\mathbb{R}}
\newcommand{\E}{\mathbb{E}}
\newcommand{\Prob}{\mathbb{P}}
\newcommand{\dd}{\mathrm{d}}
\newcommand{\T}{\mathsf{T}}
\newcommand{\Nd}{\mathcal{N}}
\newcommand{\Deltat}{\Delta_t}
\newcommand{\Sig}{\Sigma}
\newcommand{\Sigz}{\Sigma_0}
\newcommand{\Sigt}{\Sigma_t}
\newcommand{\Qz}{Q_0}
\newcommand{\Qt}{Q_t}
\newcommand{\score}{S}
\newcommand{\kl}[2]{D_{\mathrm{KL}}\!\left(#1 \,\|\, #2\right)}
\newcommand{\inner}[2]{\left\langle #1, #2 \right\rangle}
\newcommand{\norm}[1]{\left\lVert #1 \right\rVert}
